\section{Introducción}
\justifying

En la Universidad Autónoma Metropolitana Unidad Azcapotzalco, particularmente en la División de Ciencias y Artes para el Diseño (CyAD), los profesores son responsables de llenar manualmente formularios en Google Forms que luego son capturados en hojas de cálculo para gestionar documentos clave como las cartas temáticas, las autoevaluaciones docentes y los requisitos de recuperación para alumnos. Este proceso, repetitivo y propenso a errores, dificulta la organización y el seguimiento adecuado de la información académica.
Para resolver este problema, se propone el desarrollo de un sistema digital dirigido principalmente a los docentes, que automatice el llenado de datos académicos una vez que el usuario se haya identificado y seleccione el documento con el que va a trabajar. El sistema permitirá a los profesores: (1) gestionar cartas temáticas, (2) realizar su autoevaluación docente al finalizar el trimestre, y (3) subir requisitos de recuperación. Además, el sistema facilitará a los administradores el acceso a las actividades realizadas por los docentes y la generación de reportes sobre el cumplimiento de estas actividades. 

\newpage
